\section{Implementation and Evaluation}
The previous chapter outlined the design principles and methodological
foundations guiding the development of Cedar. This chapter presents the
\textbf{realized prototype}, describing its implementation,
internal architecture, and preliminary evaluation. The aim is not to
demonstrate a production-ready compiler but to validate the design's
feasibility and examine how declarative layout specifications can
operate directly within the C ecosystem. The Cedar prototype serves as a
\textit{proof of concept} for the design-science process described earlier:
it embodies a set of language features and checking mechanisms from which
insights about expressiveness, soundness, and practicality can be drawn.

This chapter begins by providing an overview of the compiler's implementation
and module organization. It then delves into descriptions of
intermediate representations and key components—layout checking, type
matching, and code generation. Finally, it reports evaluation results
obtained form layout examples and reflects on design trade-ffs and lessons
learned.

\subsection{Language Syntax and Grammar}
Cedar extends the C programming language with a small set of
\textbf{declarative constructs} for describing how data is represented
in memory. Its syntax is intentionally compact and close to C, so that
programmers can reason about physical layout without leaving 
the familiar idioms of systems programming. A Cedar source file
defines one or more \textbf{layout declaration} that specify
the mapping between a logical C type and its concrete in-memory
representation. Each layout is associated with a base C type and describes
the placement, alignment, and size of its fields using a combination
of absolute and relative directives.

\subsubsection{Core Syntax}
%Figure
This grammar captures only the layout-level constructs; standard C syntax
for type declarations is assumed to exists
in companion \texttt{.h} or \texttt{.c} files. Layout declarations may
appear in any order, and can reference previously defined layouts through
identifiers.

\subsubsection{Example}
%code snippet
This example declares a memory layout for the \texttt{Student} structure.
The layout begins at bit offset 0.
Each field specifies a C type and a placement directive:
\begin{itemize}
    \item \texttt{name} is a fixed-size character array occupying the first 256 bits.
    \item \texttt{id} is placed immediately \textit{after} \texttt{name}, inheriting its offset plus size.
    \item \texttt{grade} follows \texttt{id}.
\end{itemize}

Cedar's compiler resolves these relationships to compute
explicit bit-ranges and verify that fields neither overlap nor exceed
the enclosing structure's bounds.

\subsubsection{Design Rationale}
The syntax follows three guiding principles:
\begin{enumerate}
    \item  \textbf{Familiarity:} preserve C's lexical style so existing 
    programmers can adopt Cedar easily.
    \item \textbf{Clarity:} keep layouts readable by expressing relationships
    (\texttt{after}, \texttt{before}) rather than numeric offsets.
    \item \textbf{Analyze-ability:} restrict constructs to a small, statically
    checkable core that can be reasoned about formally in future work.
\end{enumerate}
The existing prototype currently omits recursive and parametric layouts,
focusing on straightforward records and arrays. These restrictions simplify
reasoning about offsets and enable predictable code generation. Later
chapters discuss how these constructs translate through
the compiler pipeline and
into generated C code.

%Figure on how a layout is transofrmed into intermediate representations
% and finally C code.


\subsection{Implementation Overview}
Cedar has been implemented as \textbf{modular Haskell compiler}. 
The implementation follows a multi-stage architecture. 

%Figure illustrating the flow of the compiler

\subsubsection{Structure and Tooling}
The compiler is organized into clearly separated modules:

\begin{center}
\begin{tabular}{||c c ||}
    \hline
    \textbf{Module} & \textbf{Responsibility} \\ [0.5 ex] 
    \hline\hline
    \texttt{Surface.hs} & Defines the abstract syntax for user-facing Cedar constructs. \\
    \hline
    \texttt{Lexer.x and Parser.y} & Tokenize and parse \texttt{.cedar} files. \\
    \hline
\end{tabular}
\end{center}
